\section{Content}

\subsection{Useful definitions}

In our setting, global types are automata that describe a language of MSCs. 
An \emph{arrow} is a triple $(p,q,m)\in\Procs\times\Procs\times\Messages$ with $p\ne q$;
we often write $\marrow{p}{q}{m}$ instead of $(p,q,m)$, and write $\labelalphabet$
to denote the finite set of arrows.

\begin{definition}[Global Type]
A global type $\gt$ is a deterministic finite state automaton over the alphabet $\labelalphabet$.
\end{definition}

Informally, weakly-realizable means that it's realizable and deadlock is allowed.
%% definire da linguaggio di G
\begin{definition}[Weakly-realizable]
It's the 
weak closure $L^w$ of a set $L$ of global type contains all the global type $L$ weakly implies.
The set $L$ is weakly realizable iff $L=L^w$. The global type-graph $G$ is said to be weakly
realizable if the set $L(G)$ of global type is. Thus, a weakly realizable graph already contains
all the implied scenarios.
\end{definition}


\begin{definition}[Relaxed Post Correspondence Problem (RPCP) problem]
Given $\{(v_1, w_1), (v_2, w_2), ..., (v_r, w_r)\}$, determinate whether there are indices $i_1, ..., i_m$
such that $x_{i_1}...x_{i_n}=y_{i_1}...y_{i_m}$, where $x_{i_j}, y_{i_j} \in \{v_{i_j},w_{i_j}\}$, for some
index $i_l \  x_{i_l}\neq y_{i_l}$, and for all $j\leq m,y_{i_1}...y_{i_j}$ is a prefix of $x_{i_1}...x_{i_j}$.
The relaxed problem essentially says that must have the left side of the equivalence that grows more than the right side,
we must have also some tuple that are different, and that we can choose whatever side of tuple for the composition
of the string.
\end{definition}

RPCP is proven to be an undecidable problem by reduction from the classic
PCP problem by Alur, et al. \cite{alur2005realizability}.

\subsection{Proof}

\begin{theorem}
Given a global type $G$, checking if $G$ is weakly-realizable is undecidable.
\end{theorem}

\begin{proof}
The proof proceeds via a reduction from the RPCP problem. Let's define some useful
elements for the proof.

Given a finite set $L$ of global types, where each is defined as a deterministic
finite state automaton over the alphabet $\labelalphabet$, we define $L^*$
to be the language consisting of all finite concatenations of languages in $L$.
To construct an automaton recognizing $L^*$, we introduce an initial state
$v_I$ and a terminal state $v_T$, and for each automaton in $L$, we add
a transition from $v_I$ to its initial state and from each of its final states
to $v_T$. Additionally, we connect every final state of each automaton in $L$
to the initial state of every other automaton in $L$. Thus, $L^*$ contains all
MSCs that can be obtained by concatenating MSCs from individual languages in $L$.

Given an instance $\Delta = {(v_1, w_1), \ldots, (v_m, w_m)}$ of RPCP,
we build a set $L$ of languages over 4 processes as follows.
For each pair $(v_i, w_i)$ we construct a language $\mathcal{L}_i^n$
(where $n \in {0,1}$) consisting of MSCs specified by an automaton
$A_i^n$ (described in Fig.~\ref{fig:gtype}).

In each global type $G_i^n$, process $q$ sends a message $(i,n)$ to itself,
then sends the index $i$ to process 4; process 4 then sends $(i,n)$
to process 3. After this, process 2 sends the characters of $v_i$
(or $w_i$, depending on $n$) to process 3. This communication pattern
ensures that the MSCs in $\mathcal{L}_i^n$ describe structured
interactions with explicit acknowledgments.
Observe that the communication graph of each MSC in $\mathcal{L}_i^n$
is strongly connected and involves all processes.

\begin{figure}[h]
\centering
\begin{tikzpicture}[->, node distance=24mm, on grid, auto]
  \node[state] (q0) {$q_0$};
  \node[state] (q1) [right=of q0] {$q_1$};
  \node[state] (q2) [right=of q1] {$q_2$};
  \node[state] (q3) [right=of q2] {$q_3$};
  \node[state] (q4) [right=of q3] {$q_4$};
  \node[state] (q5) [right=of q4] {$q_x$};
  \node[state] (q6) [right=of q5] {$q_y$};

  \path (q0) edge[] node[above] {\arrmess{p}{q}{(i,n)}} (q1);
  \path (q1) edge[] node[above] {\arrmess{p}{s}{i}} (q2);
  \path (q2) edge[] node[above] {\arrmess{s}{r}{(i,n)}} (q3);
  \path (q3) edge[] node[above] {\arrmess{q}{r}{x_i^1}} (q4);
  \path (q4) edge[] node[above] {\arrmess{q}{r}{...}} (q5);
  \path (q5) edge[] node[above] {\arrmess{q}{r}{x_i^c}} (q6);
\end{tikzpicture}
\caption{The global type $G^n_i$ where $n\in\{0,1\}$ and $x\in\{v,w\}$.}
\label{fig:gtype}
\end{figure}

We need to prove:
\begin{center}
$\Delta \in \text{RPCP}$ if and only if the language $L^*$ is not realizable.
\end{center}

\begin{itemize}

  \item[$\Rightarrow$] Suppose there is a solution to RPCP: $\Delta = \{i_1,\ldots,i_k$\}
  with $v_{i_1}\ldots v_{i_k} = w_{i_1}\ldots w_{i_k}$. Then, ...
  
  \item[$\Leftarrow$] If the language $L^*$ is not realizable,...

\end{itemize}

\textbf{Lemma.} Each $\mathcal{L}_i^n$ is a synch-language.

\begin{proof}
We show that the MSCs in $\mathcal{L}_i^n$ admit only synchronous linearizations.
Due to the presence of ACK messages, send and receive
events are closely paired. The only possible reordering involves
the ACK from \arrmess{p}{s}{i}, ...
\end{proof}


\end{proof}


